\chapter{吉田第六版目录与版本对照}\label{app:yosida-map}
\section{主线章节与参考版本}
本书第一部分保留中译本第零章至第十四章的原章次。以下页码为原书印刷页码；本地两个 PDF 的正文第 1 页均位于 PDF 第 16 页，因此正文部分的 PDF 页序等于印刷页码加 15。前置页不使用这一换算。中译本目录位于 PDF 第 10--15 页，英文版目录位于 PDF 第 11--15 页。
\begin{center}\small
\begin{tabular}{p{0.08\linewidth}p{0.53\linewidth}p{0.14\linewidth}p{0.14\linewidth}}
\toprule 章 & 主题 & 中译本页码 & 英文版页码\\\midrule
0 & 预备知识 & 1--22 & 1--22 \\
1 & 半范数 & 23--70 & 23--67 \\
2 & Baire--Hausdorff 定理的应用 & 71--84 & 68--80 \\
3 & 正交投影及 F. Riesz 表示定理 & 85--106 & 81--101 \\
4 & Hahn--Banach 定理 & 107--124 & 102--119 \\
5 & 强收敛和弱收敛 & 125--150 & 119--145 \\
6 & Fourier 变换和微分方程 & 151--200 & 145--192 \\
7 & 对偶算子 & 201--218 & 193--208 \\
8 & 预解式和谱 & 219--242 & 209--231 \\
9 & 半群的分析理论 & 243--286 & 231--274 \\
10 & 紧算子 & 287--306 & 274--293 \\
11 & 赋范环和谱表示 & 307--376 & 294--362 \\
12 & 线性空间中其他表示定理 & 377--396 & 362--378 \\
13 & 遍历理论和扩散理论 & 397--436 & 379--418 \\
14 & 发展方程的积分 & 437--486 & 418--465 \\
\bottomrule\end{tabular}\end{center}
\section{目录核对说明}
章节名与各节主题以中译本为准，统一连字符、标点与少数人名拼写的排版。不使用英文版页码替代中译本页码。第五章附录的四个子目在中译本总目录中未展开，依据英文版目录补足，已在第五章附录中保留。空间记号表依据英文版书后目录建立辅助占位。
英文章名及页码用于辅助核对目录，不替代中译本页码。正文复查确认几处共页：英文第 119 页上部为第四章结尾、下部开始第五章；第 145 页由第五章接到第六章；第 231 页由第八章接到第九章；第 274 页由第九章接到第十章。本次又确认第 362 页上部仍为第十一章 Littlewood 推论证明结尾，下部才开始第十二章，故第十一章英文范围修正为 294--362 页，第十二章仍从 362 页起。本轮确认英文第 418 页上部仍为第十三章抽象位势说明、下部开始第十四章，故第十三章改为 379--418 页，第十四章起始 418 不变。这些范围允许重叠，编号节起始页未改。

目前第零至第十三章共 125/132 个编号节已有首轮正文，配 218 道附解答习题，第五、十章的章内附录亦已写入。第十一章涵盖交换赋范代数、谱测度、Friedrichs 扩张、紧群表示及 Wiener 稠密定理，配三十二道习题。第十二章六节及十四道习题连接凸集端点、序收敛与格的表示。第十三章九节与十八道习题连接遍历、扩散与位势。各章的实际查阅页、数学自查、外部前提与尚未展开的证明均分别记录于对应 COVERAGE，不把已有正文或编译通过等同于独立数学审校。第十四章的七节仍保留占位。

\section{广义函数专题的接入位置}
第二部分为第十五章至第二十二章，沿用现有八个专题文件。第一部分中 I.8、I.12--14、II.3、III.11、VI.1--3 等主题仍保留其原书位置；相关节用标签与对应专题章衔接。这使主线目录保持完整，也便于后续逐节写入吉田原书的特定证明与论述。
\section{占位完成后的更新规则}
正文写成后删除占位提示，把页码转入出处说明；保留章、节标题及标签。新增定理与习题使用独立标签，章号随全书顺序自动更新。
